\section{代表性攻击路径一：Prompt Injection 到真实执行}

本讲对 prompt injection 做了系统升级，不再局限 ``模型被诱导说错话''，而是分析 ``模型被诱导执行错动作''。尤其是 indirect prompt injection：攻击载荷藏在网页、邮件、文档或 issue 中，Agent 在合法读取后将恶意指令当作环境信息执行。

\begin{importantbox}{Indirect Prompt Injection 的危险性}
攻击者无需突破系统边界，只需投递 ``被 Agent 合法读取'' 的恶意内容，就可能触发后续工具调用链。这使得攻击成本低、隐蔽性高、扩散速度快。
\end{importantbox}

课程给出的 web agent 示例中，Agent 被诱导忽略用户原始意图，转而访问恶意站点并下载恶意内容。这里的关键是 ``攻击链'' 而非单点漏洞：
\begin{enumerate}
    \item 污染输入源（网页内容、issue 文本、邮件正文）；
    \item 触发计划偏移（从任务目标转向攻击者目标）；
    \item 借助高权限工具执行危险动作；
    \item 利用持久状态扩大后续危害。
\end{enumerate}

\begin{warningbox}{工程中最容易忽略的一点}
很多团队只过滤 ``用户输入''，却忽视 ``环境输入''。实际上，环境输入（网页、检索结果、工具返回）才是 indirect injection 的主战场。
\end{warningbox}

\subsection{本章小结}
Prompt injection 在 Agent 时代的本质是执行链劫持。防御重点应从文本过滤扩展到任务状态约束与工具执行约束。

